\chapter{Background}

\note{
TODO: This is the background section from PACT paper... need look over and revise.
}

\subsection{Sparse Matrix Storage Formats}
The four formats considered in this work represent the standard vocabulary of sparse
matrix storage. Compressed Sparse Row (CSR) stores non-zeros row by row, with a row
pointer array enabling efficient traversal of row entries.
Compressed Sparse Column (CSC) is the column-major analog: columns are stored
contiguously, making it efficient for column-wise access patterns. Coordinate format
(COO) stores each non-zero as an explicit $(\text{row}, \text{col}, \text{value})$
triple, meaning pointer compression but consuming more memory per entry and
imposing a scan cost proportional to the total non-zero count during row-wise access.
Dense format stores all elements explicitly in a flat row-major array, eliminating
indexing overhead entirely and exposing maximum memory locality, making it optimal
when most entries are non-zero, unsurprisingly. Each format encodes the same mathematical
object but makes different trade-offs between storage overhead,
access pattern regularity, and the cost of skipping structural zeros.
The standard reference for these formats in the context of sparse linear algebra
compilers is \cite{taco_formats}, who provide a unified description of
format capabilities in terms of per-dimension storage modes.

\subsection{SpGEMM}
SpGEMM computes $C = A \times B$ where $A$ and $B$ are sparse matrices. Unlike sparse
matrix-vector multiplication (SpMV), SpGEMM involves two sparse operands whose
non-zero structure interacts during multiplication: the output sparsity pattern is not
known in advance and must be computed symbolically before numeric accumulation can
begin. The operation appears as a core subroutine in algebraic multigrid solvers,
triangle counting, multi-source breadth-first search, and high-dimensional similarity
search, among many other applications~\cite{spgemm_survey}.
Buluç and Gilbert~\cite{buluc_gilbert}
provide foundational parallel algorithms for SpGEMM.
The most performance-relevant observation for this work is that SpGEMM is overwhelmingly
memory-bound on current hardware: the irregular, format-dependent access patterns of
sparse operands generate large numbers of short, unpredictable memory requests, keeping
arithmetic units underutilized.

\subsection{BLIS}
The BLAS-like Library Instantiation Software (BLIS) framework~\cite{blis} provides a
portable infrastructure for implementing dense level-3 BLAS operations. Its key insight
is that virtually all computation in dense matrix multiplication can be expressed in
terms of a small, hardware-specific micro-kernel operating on dense register-sized
panels, surrounded by a fixed five-loop blocking structure parameterized by cache-level
constants $m_c$, $n_c$, $k_c$ and micro-kernel block sizes $m_r$ and $n_r$. This
separation allows the bulk of the framework to be written once in portable C while
performance-critical sections are tuned per architecture. Our implementation preserves
this five-loop structure and blocking hierarchy, replacing only the innermost
micro-kernel with a dispatch over heterogeneous sparse formats.

\subsection{Roofline Model}
The roofline model~\cite{roofline} provides an upper bound on achievable kernel
performance as a function of arithmetic intensity (FLOPs per byte of memory traffic).
For a given kernel and machine, the model draws a ``roofline'' as the minimum of the
peak compute throughput and the product of memory bandwidth and arithmetic intensity. A
kernel is compute-bound if its arithmetic intensity exceeds the machine's
compute-to-bandwidth ratio, and memory-bound otherwise. SpGEMM kernels consistently
fall in the memory-bound regime due to low data reuse and irregular access patterns, a
finding confirmed by our results in Section 4.

\subsection{TACO}
The Tensor Algebra Compiler (TACO)~\cite{taco} automatically generates efficient
kernels for any compound tensor algebra expression over mixed dense and sparse tensors,
including all format combinations relevant to this work. TACO provides a principled
formalization of how iteration spaces and index expressions vary across storage formats,
and its generated kernels served as the starting point and correctness reference for our
hand-instrumented implementations.

\subsection{Kronecker Graphs}
Kronecker graphs~\cite{kronecker} are a family of self-similar synthetic networks
generated by repeated Kronecker product of a small initiator matrix with itself. They
were introduced as a model for real-world networks because their recursive structure
naturally produces the heavy-tailed degree distributions and small diameters observed
in internet topology graphs and social networks. The AS-Newman autonomous systems
dataset used in our experiments is a standard benchmark graph commonly represented as
a Kronecker graph instance.

\subsection{Optimization}
Format selection in this work is formulated as a Mixed Integer Nonlinear Program
(MINLP) and solved using Pyomo~\cite{pyomo} with MindtPy~\cite{mindtpy} as the MINLP
solver back-end, using GLPK for MIP subproblems and IPOPT for the NLP relaxations.
